% Ch2 WS3 -- extra nomenclature practice, built on Fig 2.21, Tables 2.5,
% 2.7, 2.8. Supplements ws2: the reference tables are built by the student
% here rather than supplied, and the item set is larger (94) and disjoint
% from ws2's wherever a different compound would do.
\documentclass{apchem-worksheet}

\apcunitword{Chapter}
\apcunit{2}
\apcunittitle{Atoms, Molecules, and Ions}
\apcdoctitle{Worksheet 3 \textbullet\ Naming Simple Compounds --- Extra Practice}
\apcsubtitle{Zumdahl \S2.8 \textbullet\ build the tables, then 94 items}

\begin{document}
\wsheader[Do Part 1 from memory first, then check it against Fig.~2.21 and
Table~2.5 in the text. Everything in Parts 3--4 is faster once Part 1 is
automatic.]

\begin{apcnote}[How to use this sheet]
Naming is not a reasoning skill, it is a \textbf{recall} skill sitting on
top of one decision. The decision is always the same:

\begin{center}
\renewcommand{\arraystretch}{1.2}
\begin{tabular}{@{}llp{7.4cm}@{}}
\hline
\textbf{Type} & \textbf{What it is} & \textbf{How you know} \\
\hline
I & ionic, fixed-charge metal & metal in group 1A, 2A, 3A, or \ce{Ag},
  \ce{Zn}, \ce{Cd} --- no Roman numeral \\
II & ionic, variable-charge metal & transition metal (or \ce{Sn},
  \ce{Pb}) --- Roman numeral \emph{required} \\
III & binary covalent & two nonmetals --- prefixes, and only here \\
A & acid & formula begins with \ce{H} and is aqueous \\
\hline
\end{tabular}
\end{center}

Almost every naming error is a \textbf{type} error, not a memory error ---
prefixes used on an ionic compound, or a Roman numeral left off a
transition metal.
\end{apcnote}

\section*{\sffamily Part 1 \textbullet\ Build the reference tables}

\problem[]{\textbf{Common cations} (Fig.~2.21). Fill in the charge each
element forms. The left block is fixed-charge; the right block is
variable-charge and needs a Roman numeral every time.}

\renewcommand{\arraystretch}{1.35}
\noindent\begin{tabular}{@{}p{3.4cm}p{3.6cm}p{3.4cm}p{4.2cm}@{}}
\toprule
\multicolumn{2}{@{}l}{\textbf{Fixed charge (Type I)}} &
\multicolumn{2}{l@{}}{\textbf{Variable charge (Type II)}}\\
\midrule
Group 1A (\ce{Li}, \ce{Na}, \ce{K}, \ce{Rb}, \ce{Cs}) &
  \answerblank[2.2cm]{$1+$} &
  \ce{Cr} & \answerblank[3.2cm]{$2+$, $3+$}\\
Group 2A (\ce{Mg}, \ce{Ca}, \ce{Sr}, \ce{Ba}) &
  \answerblank[2.2cm]{$2+$} &
  \ce{Mn} & \answerblank[3.2cm]{$2+$, $3+$}\\
Group 3A (\ce{Al}, \ce{Ga}) &
  \answerblank[2.2cm]{$3+$} &
  \ce{Fe} & \answerblank[3.2cm]{$2+$, $3+$}\\
\ce{Ag} & \answerblank[2.2cm]{$1+$} &
  \ce{Co} & \answerblank[3.2cm]{$2+$, $3+$}\\
\ce{Zn} & \answerblank[2.2cm]{$2+$} &
  \ce{Cu} & \answerblank[3.2cm]{$1+$, $2+$}\\
\ce{Cd} & \answerblank[2.2cm]{$2+$} &
  \ce{Sn}, \ce{Pb} & \answerblank[3.2cm]{$2+$, $4+$}\\
 & & \ce{Hg} & \answerblank[3.6cm]{\ce{Hg2^2+}, \ce{Hg^2+}}\\
\bottomrule
\end{tabular}

\vspace{0.3em}
\noindent\textbf{Common monatomic anions.} Give the ion and its name:

\begin{parts}
  \item \ce{N}: \answerblank[4.6cm]{\ce{N^3-}, nitride} \quad
        \ce{O}: \answerblank[4.6cm]{\ce{O^2-}, oxide}
  \item \ce{S}: \answerblank[4.6cm]{\ce{S^2-}, sulfide} \quad
        \ce{F}: \answerblank[4.6cm]{\ce{F-}, fluoride}
  \item \ce{Cl}: \answerblank[4.6cm]{\ce{Cl-}, chloride} \quad
        \ce{Br}: \answerblank[4.6cm]{\ce{Br-}, bromide}
  \item \ce{I}: \answerblank[4.6cm]{\ce{I-}, iodide} \quad
        rule: \answerblank[5.4cm]{element root $+$ \emph{-ide}}
\end{parts}

\problem[]{\textbf{Polyatomic ions} (Table~2.5). Write the name beside each
formula. These are the ones worth knowing cold.}

\renewcommand{\arraystretch}{1.3}
\noindent\begin{tabular}{@{}p{1.9cm}p{5.2cm}p{1.9cm}p{5.2cm}@{}}
\toprule
\textbf{Ion} & \textbf{Name} & \textbf{Ion} & \textbf{Name}\\
\midrule
\ce{NH4+} & \answerblank[4.4cm]{ammonium} &
  \ce{CO3^2-} & \answerblank[4.4cm]{carbonate}\\
\ce{NO2-} & \answerblank[4.4cm]{nitrite} &
  \ce{HCO3-} & \answerblank[4.4cm]{hydrogen carbonate}\\
\ce{NO3-} & \answerblank[4.4cm]{nitrate} &
  \ce{ClO-} & \answerblank[4.4cm]{hypochlorite}\\
\ce{SO3^2-} & \answerblank[4.4cm]{sulfite} &
  \ce{ClO2-} & \answerblank[4.4cm]{chlorite}\\
\ce{SO4^2-} & \answerblank[4.4cm]{sulfate} &
  \ce{ClO3-} & \answerblank[4.4cm]{chlorate}\\
\ce{HSO4-} & \answerblank[4.4cm]{hydrogen sulfate} &
  \ce{ClO4-} & \answerblank[4.4cm]{perchlorate}\\
\ce{OH-} & \answerblank[4.4cm]{hydroxide} &
  \ce{C2H3O2-} & \answerblank[4.4cm]{acetate}\\
\ce{CN-} & \answerblank[4.4cm]{cyanide} &
  \ce{MnO4-} & \answerblank[4.4cm]{permanganate}\\
\ce{PO4^3-} & \answerblank[4.4cm]{phosphate} &
  \ce{Cr2O7^2-} & \answerblank[4.4cm]{dichromate}\\
\ce{HPO4^2-} & \answerblank[4.4cm]{hydrogen phosphate} &
  \ce{CrO4^2-} & \answerblank[4.4cm]{chromate}\\
\ce{H2PO4-} & \answerblank[4.4cm]{dihydrogen phosphate} &
  \ce{O2^2-} & \answerblank[4.4cm]{peroxide}\\
\ce{SCN-} & \answerblank[4.4cm]{thiocyanate} &
  \ce{C2O4^2-} & \answerblank[4.4cm]{oxalate}\\
\ce{Hg2^2+} & \answerblank[4.4cm]{mercury(I)} &
  \ce{S2O3^2-} & \answerblank[4.4cm]{thiosulfate}\\
\bottomrule
\end{tabular}

\begin{apcnote}[The oxyanion pattern --- learn this instead of memorizing]
Where a series of anions differ only in oxygen count, the names follow one
pattern. Two members: fewer oxygens $\to$
\answerblank[1.6cm]{\emph{-ite}}, more $\to$
\answerblank[1.8cm]{\emph{-ate}}. More than two: the extremes take
\answerblank[2cm]{\emph{hypo-}} (fewest) and
\answerblank[1.6cm]{\emph{per-}} (most).

The chlorine series is the one to hold in memory, because every other
halogen copies it exactly:
\[ \underbrace{\ce{ClO-}}_{\text{hypochlorite}} \quad
   \underbrace{\ce{ClO2-}}_{\text{chlorite}} \quad
   \underbrace{\ce{ClO3-}}_{\text{chlorate}} \quad
   \underbrace{\ce{ClO4-}}_{\text{perchlorate}} \]
Know \ce{ClO3-} is chlorate and \ce{SO4^2-} is sulfate, and the rest of
each series follows without memorizing anything else.
\end{apcnote}

\section*{\sffamily Part 2 \textbullet\ Acids}

An acid is one or more \ce{H+} ions attached to an anion, so \textbf{the
acid name is decided entirely by the anion's name}:

\begin{center}
\renewcommand{\arraystretch}{1.25}
\begin{tabular}{@{}lll@{}}
\toprule
\textbf{Anion ending} & \textbf{Acid name} & \textbf{Example}\\
\midrule
\emph{-ide} (no oxygen) & \textbf{hydro}(root)\textbf{ic} acid &
  chloride $\to$ hydrochloric acid\\
\emph{-ite} & (root)\textbf{ous} acid & sulfite $\to$ sulfurous acid\\
\emph{-ate} & (root)\textbf{ic} acid & sulfate $\to$ sulfuric acid\\
\bottomrule
\end{tabular}
\end{center}

\problem[]{\textbf{Acids without oxygen} (Table~2.7). Complete every blank
column.}

\renewcommand{\arraystretch}{1.3}
\noindent\begin{tabular}{@{}p{2.4cm}p{4.6cm}p{6.4cm}@{}}
\toprule
\textbf{Acid} & \textbf{Anion (name)} & \textbf{Acid name}\\
\midrule
\ce{HF} & \answerblank[3.6cm]{fluoride} &
  \answerblank[5cm]{hydrofluoric acid}\\
\ce{HCl} & \answerblank[3.6cm]{chloride} &
  \answerblank[5cm]{hydrochloric acid}\\
\ce{HBr} & \answerblank[3.6cm]{bromide} &
  \answerblank[5cm]{hydrobromic acid}\\
\ce{HI} & \answerblank[3.6cm]{iodide} &
  \answerblank[5cm]{hydroiodic acid}\\
\ce{HCN} & \answerblank[3.6cm]{cyanide} &
  \answerblank[5cm]{hydrocyanic acid}\\
\ce{H2S} & \answerblank[3.6cm]{sulfide} &
  \answerblank[5cm]{hydrosulfuric acid}\\
\bottomrule
\end{tabular}

{\sffamily\footnotesize These names apply to the \emph{aqueous solutions}.
Gaseous \ce{HCl} is hydrogen chloride; \ce{HCl(aq)} is hydrochloric acid.}

\problem[]{\textbf{Oxygen-containing acids} (Table~2.8, extended with the
chlorine series). Complete every blank column.}

\renewcommand{\arraystretch}{1.3}
\noindent\begin{tabular}{@{}p{2.4cm}p{4.6cm}p{6.4cm}@{}}
\toprule
\textbf{Acid} & \textbf{Anion (name)} & \textbf{Acid name}\\
\midrule
\ce{HNO3} & \answerblank[3.6cm]{nitrate} &
  \answerblank[5cm]{nitric acid}\\
\ce{HNO2} & \answerblank[3.6cm]{nitrite} &
  \answerblank[5cm]{nitrous acid}\\
\ce{H2SO4} & \answerblank[3.6cm]{sulfate} &
  \answerblank[5cm]{sulfuric acid}\\
\ce{H2SO3} & \answerblank[3.6cm]{sulfite} &
  \answerblank[5cm]{sulfurous acid}\\
\ce{H3PO4} & \answerblank[3.6cm]{phosphate} &
  \answerblank[5cm]{phosphoric acid}\\
\ce{HC2H3O2} & \answerblank[3.6cm]{acetate} &
  \answerblank[5cm]{acetic acid}\\
\ce{HClO4} & \answerblank[3.6cm]{perchlorate} &
  \answerblank[5cm]{perchloric acid}\\
\ce{HClO3} & \answerblank[3.6cm]{chlorate} &
  \answerblank[5cm]{chloric acid}\\
\ce{HClO2} & \answerblank[3.6cm]{chlorite} &
  \answerblank[5cm]{chlorous acid}\\
\ce{HClO} & \answerblank[3.6cm]{hypochlorite} &
  \answerblank[5cm]{hypochlorous acid}\\
\bottomrule
\end{tabular}

\begin{apctrap}
\textbf{``I ate it, so it was ic.''} The \emph{-ate} anion gives the
\emph{-ic} acid. The commonest slip on this topic is naming \ce{H2SO3}
sulfuric acid --- it contains sulf\emph{ite}, so it is sulfur\emph{ous}
acid.
\end{apctrap}

\section*{\sffamily Part 3 \textbullet\ Worked examples}

\begin{workedexample}[Worked example 1: \ce{Fe(ClO3)2}, one step at a time]
\textbf{1. Type?} \ce{Fe} is a transition metal $\Rightarrow$ Type II, so a
Roman numeral is required and no prefixes are used.

\textbf{2. Anion?} \ce{ClO3-} is chlorate, charge $1-$. There are two of
them, so the negative charge totals $2-$.

\textbf{3. Cation charge, by balance?} \ce{Fe} must be $2+$.

\textbf{4. Assemble.} \textbf{Iron(II) chlorate.}

\emph{Note what supplied the Roman numeral} --- not the formula's
subscript on iron, but the total charge on the anions. The subscript 2
belongs to \ce{ClO3-}, not to \ce{Fe}.
\end{workedexample}

\begin{workedexample}[Worked example 2: reverse --- ammonium sulfate]
\textbf{1. Ions.} Ammonium is \ce{NH4+}; sulfate is \ce{SO4^2-}.

\textbf{2. Balance.} One $2-$ needs two $1+$, so two ammonium ions.

\textbf{3. Assemble.} A polyatomic ion taken more than once needs
parentheses: \ce{(NH4)2SO4}.

Writing \ce{NH42SO4} would read as forty-two hydrogens. Parentheses are not
decoration.
\end{workedexample}

\begin{workedexample}[Worked example 3: three lookalikes]
\[ \ce{Na2SO3} \qquad \ce{Na2SO4} \qquad \ce{Na2S2O3} \]
Sodium is fixed-charge in all three, so the whole answer sits in the anion:
sulf\emph{ite}, sulf\emph{ate}, and \emph{thio}sulfate respectively ---
\textbf{sodium sulfite}, \textbf{sodium sulfate},
\textbf{sodium thiosulfate}.

The prefix \emph{thio-} means an oxygen has been replaced by a sulfur:
compare \ce{SO4^2-} with \ce{S2O3^2-}.
\end{workedexample}

\section*{\sffamily Part 4 \textbullet\ 94 items}

\problem[]{\textbf{Type identification.} Label each I, II, III, or A. Do
not name them yet.}
\begin{parts}
  \item \ce{K2S} \answerblank[2cm]{I}
  \item \ce{CoF3} \answerblank[2cm]{II}
  \item \ce{N2O5} \answerblank[2cm]{III}
  \item \ce{HClO4(aq)} \answerblank[2cm]{A}
  \item \ce{Ba(NO3)2} \answerblank[2cm]{I}
  \item \ce{SnBr4} \answerblank[2cm]{II}
\end{parts}

\problem[]{\textbf{Name each (Type I).}}
\begin{parts}
  \item \ce{K2S}: \answerblank[5.5cm]{potassium sulfide}
  \item \ce{SrI2}: \answerblank[5.5cm]{strontium iodide}
  \item \ce{Li3N}: \answerblank[5.5cm]{lithium nitride}
  \item \ce{AlBr3}: \answerblank[5.5cm]{aluminum bromide}
  \item \ce{Ag2S}: \answerblank[5.5cm]{silver sulfide}
\end{parts}

\problem[]{\textbf{Name each (Type II).} Get the metal's charge from the
anions first.}
\begin{parts}
  \item \ce{CoF3}: \answerblank[5.5cm]{cobalt(III) fluoride}
  \item \ce{Cu2S}: \answerblank[5.5cm]{copper(I) sulfide}
  \item \ce{PbO2}: \answerblank[5.5cm]{lead(IV) oxide}
  \item \ce{Fe2O3}: \answerblank[5.5cm]{iron(III) oxide}
  \item \ce{SnBr4}: \answerblank[5.5cm]{tin(IV) bromide}
  \item \ce{Hg2Cl2}: \answerblank[5.5cm]{mercury(I) chloride}
\end{parts}

\problem[]{\textbf{Name each compound containing a polyatomic ion.}}
\begin{parts}
  \item \ce{Ba(NO3)2}: \answerblank[6.4cm]{barium nitrate}
  \item \ce{Na2Cr2O7}: \answerblank[6.4cm]{sodium dichromate}
  \item \ce{Fe(C2H3O2)2}: \answerblank[6.4cm]{iron(II) acetate}
  \item \ce{(NH4)2SO3}: \answerblank[6.4cm]{ammonium sulfite}
  \item \ce{Ca(H2PO4)2}: \answerblank[6.4cm]{calcium dihydrogen phosphate}
  \item \ce{KMnO4}: \answerblank[6.4cm]{potassium permanganate}
\end{parts}

\problem[]{\textbf{Name each binary covalent compound (Type III).}}
\begin{parts}
  \item \ce{N2O5}: \answerblank[5.5cm]{dinitrogen pentoxide}
  \item \ce{SF6}: \answerblank[5.5cm]{sulfur hexafluoride}
  \item \ce{PBr3}: \answerblank[5.5cm]{phosphorus tribromide}
  \item \ce{CS2}: \answerblank[5.5cm]{carbon disulfide}
  \item \ce{Cl2O7}: \answerblank[5.5cm]{dichlorine heptoxide}
\end{parts}

\problem[]{\textbf{Name each acid} (all aqueous).}
\begin{parts}
  \item \ce{HI(aq)}: \answerblank[5.5cm]{hydroiodic acid}
  \item \ce{HClO2(aq)}: \answerblank[5.5cm]{chlorous acid}
  \item \ce{H2CO3(aq)}: \answerblank[5.5cm]{carbonic acid}
  \item \ce{HC2H3O2(aq)}: \answerblank[5.5cm]{acetic acid}
  \item \ce{HNO2(aq)}: \answerblank[5.5cm]{nitrous acid}
  \item \ce{H2S(aq)}: \answerblank[5.5cm]{hydrosulfuric acid}
\end{parts}

\problem[]{\textbf{Write the formula} for each name.}
\begin{parts}
  \item cobalt(II) hydroxide: \answerblank[4cm]{\ce{Co(OH)2}}
  \item ammonium dichromate: \answerblank[4cm]{\ce{(NH4)2Cr2O7}}
  \item tetraphosphorus decoxide: \answerblank[4cm]{\ce{P4O10}}
  \item hydrocyanic acid: \answerblank[4cm]{\ce{HCN(aq)}}
  \item magnesium perchlorate: \answerblank[4cm]{\ce{Mg(ClO4)2}}
  \item iron(III) sulfate: \answerblank[4cm]{\ce{Fe2(SO4)3}}
  \item sulfur dioxide: \answerblank[4cm]{\ce{SO2}}
  \item lead(II) chromate: \answerblank[4cm]{\ce{PbCrO4}}
  \item perchloric acid: \answerblank[4cm]{\ce{HClO4(aq)}}
  \item sodium thiosulfate: \answerblank[4cm]{\ce{Na2S2O3}}
\end{parts}

\problem[]{\textbf{Extend the pattern.} Bromine forms the same oxyanion
series as chlorine.}
\begin{parts}
  \item Name \ce{BrO-}, \ce{BrO2-}, \ce{BrO3-}, \ce{BrO4-} in order.
        \answerlines[2]{Hypobromite, bromite, bromate, perbromate --- the
        chlorine pattern transferred without change.}
  \item Name the acids \ce{HBrO}, \ce{HBrO2}, \ce{HBrO3}, \ce{HBrO4}.
        \answerlines[2]{Hypobromous, bromous, bromic, and perbromic acid
        --- \emph{-ite} gives \emph{-ous} and \emph{-ate} gives
        \emph{-ic}, with the prefixes carried over.}
\end{parts}

\problem[]{\textbf{Error hunt.} Each name is wrong. Give the correct name
and say which rule was broken.}
\begin{parts}
  \item \ce{CuCl2}, ``copper chloride''
        \answerlines[2]{Copper(II) chloride --- copper is a variable-charge
        metal, so the Roman numeral is required.}
  \item \ce{Al2O3}, ``dialuminum trioxide''
        \answerlines[2]{Aluminum oxide --- prefixes are used only for
        binary covalent compounds; ionic formulas follow from charge
        balance.}
  \item \ce{Ca(NO2)2}, ``calcium nitrate''
        \answerlines[2]{Calcium nitrite --- \ce{NO2-} is nitrite;
        \ce{NO3-} is nitrate.}
  \item \ce{H2SO3(aq)}, ``sulfuric acid''
        \answerlines[2]{Sulfurous acid --- the anion is sulfite, and an
        \emph{-ite} anion gives an \emph{-ous} acid.}
\end{parts}

\problem[]{\textbf{Lookalikes.} Name \emph{both} members of each pair. One
subscript is the entire difference.}
\begin{parts}
  \item \ce{CuCl} / \ce{CuCl2}
        \answerblank[8.4cm]{copper(I) chloride / copper(II) chloride}
  \item \ce{FeO} / \ce{Fe2O3}
        \answerblank[8.4cm]{iron(II) oxide / iron(III) oxide}
  \item \ce{Na2CO3} / \ce{NaHCO3}
        \answerblank[9.6cm]{sodium carbonate / sodium hydrogen carbonate}
  \item \ce{Na2SO3} / \ce{Na2SO4}
        \answerblank[8.4cm]{sodium sulfite / sodium sulfate}
\end{parts}

\problem[]{\textbf{Mixed list.} Types are deliberately jumbled --- decide
the type before naming. This is the format the exam uses.}
\begin{parts}
  \item \ce{ICl}: \answerblank[6.4cm]{iodine monochloride}
  \item \ce{NH4NO2}: \answerblank[6.4cm]{ammonium nitrite}
  \item \ce{Co2S3}: \answerblank[6.4cm]{cobalt(III) sulfide}
  \item \ce{Pb3(PO4)2}: \answerblank[6.4cm]{lead(II) phosphate}
  \item \ce{HC2H3O2(aq)}: \answerblank[6.4cm]{acetic acid}
  \item \ce{Zn(OH)2}: \answerblank[6.4cm]{zinc hydroxide}
  \item \ce{P4O6}: \answerblank[6.4cm]{tetraphosphorus hexoxide}
  \item \ce{CrCl3}: \answerblank[6.4cm]{chromium(III) chloride}
\end{parts}

\problem[]{\textbf{Transfer the pattern again.} Elements in the same family
form oxyanions of the same general formula, named the same way.}
\begin{parts}
  \item \ce{SeO4^2-}: \answerblank[4cm]{selenate} \quad
        \ce{SeO3^2-}: \answerblank[4cm]{selenite}
  \item \ce{TeO4^2-}: \answerblank[4cm]{tellurate} \quad
        \ce{TeO3^2-}: \answerblank[4cm]{tellurite}
  \item Which chlorine oxyanion did you use as the model, and why does it
        work for selenium? \answerlines[2]{Chlorate/chlorite --- the
        \emph{-ate} form is the reference and the one with one fewer
        oxygen takes \emph{-ite}. The suffix encodes relative oxygen
        count, not the element, so it transfers to any such series.}
  \item Name \ce{IO-}, \ce{IO2-}, \ce{IO3-}, \ce{IO4-} in order.
        \answerlines[2]{Hypoiodite, iodite, iodate, periodate.}
  \item Name the acids \ce{HIO}, \ce{HIO2}, \ce{HIO3}, \ce{HIO4}.
        \answerlines[2]{Hypoiodous, iodous, iodic, and periodic acid.}
  \item \ce{SeO4^2-} carries a $2-$ charge while \ce{ClO4-} carries $1-$.
        Does that change how either is named?
        \answerblank[8.3cm]{no --- the name tracks oxygen count only}
\end{parts}

\problem[]{\textbf{Write the formula} for each name (second set).}
\begin{parts}
  \item potassium hydrogen sulfate: \answerblank[4cm]{\ce{KHSO4}}
  \item barium peroxide: \answerblank[4cm]{\ce{BaO2}}
  \item nickel(II) cyanide: \answerblank[4cm]{\ce{Ni(CN)2}}
  \item dinitrogen tetroxide: \answerblank[4cm]{\ce{N2O4}}
  \item chromium(III) carbonate: \answerblank[4cm]{\ce{Cr2(CO3)3}}
  \item sodium dihydrogen phosphate: \answerblank[4cm]{\ce{NaH2PO4}}
  \item hydrosulfuric acid: \answerblank[4cm]{\ce{H2S(aq)}}
  \item silicon tetrachloride: \answerblank[4cm]{\ce{SiCl4}}
  \item mercury(II) oxide: \answerblank[4cm]{\ce{HgO}}
  \item ammonium oxalate: \answerblank[4cm]{\ce{(NH4)2C2O4}}
\end{parts}

\problem[]{\textbf{Acid from its anion.} Each aqueous acid below is one or
more \ce{H+} attached to the anion named. Give the acid's formula and its
name.}
\begin{parts}
  \item \ce{PO4^3-}: \answerblank[7.4cm]{\ce{H3PO4}, phosphoric acid}
  \item \ce{ClO2-}: \answerblank[7.4cm]{\ce{HClO2}, chlorous acid}
  \item \ce{CN-}: \answerblank[7.4cm]{\ce{HCN}, hydrocyanic acid}
  \item \ce{C2O4^2-}: \answerblank[7.4cm]{\ce{H2C2O4}, oxalic acid}
  \item \ce{S2O3^2-}: \answerblank[7.4cm]{\ce{H2S2O3}, thiosulfuric acid}
  \item How many \ce{H+} does the anion's charge tell you to attach?
        \answerblank[6.4cm]{as many as the negative charge}
\end{parts}

\problem[]{\textbf{Error hunt, second set.} Give the correct name and the
rule that was broken.}
\begin{parts}
  \item \ce{SO3}, ``sulfur(VI) oxide''
        \answerlines[2]{Sulfur trioxide --- both elements are nonmetals, so
        this is covalent: prefixes are used and Roman numerals are not.}
  \item \ce{K2O}, ``potassium(I) oxide''
        \answerlines[2]{Potassium oxide --- potassium is a fixed-charge
        metal, so no Roman numeral is written.}
  \item \ce{Ca(ClO3)2}, ``calcium chlorite''
        \answerlines[2]{Calcium chlorate --- \ce{ClO3-} is chlorate;
        \ce{ClO2-} is chlorite.}
  \item \ce{SnO}, ``tin oxide''
        \answerlines[2]{Tin(II) oxide --- tin is variable-charge, and one
        \ce{O^2-} forces \ce{Sn^2+}.}
  \item \ce{Li2O2}, ``lithium oxide''
        \answerlines[2]{Lithium peroxide --- the anion is \ce{O2^2-}
        (peroxide), not \ce{O^2-} (oxide).}
  \item \ce{H2SO4(aq)}, ``sulfate acid''
        \answerlines[2]{Sulfuric acid --- an \emph{-ate} anion gives an
        \emph{-ic} acid; the anion name is never used unchanged.}
\end{parts}

\problem[]{\textbf{Say why.} One or two sentences each.}
\begin{parts}
  \item Why does \ce{ZnCl2} need no Roman numeral when \ce{FeCl2} does?
        \answerlines[2]{Zinc forms only \ce{Zn^2+}, so the charge is not in
        doubt. Iron forms \ce{Fe^2+} and \ce{Fe^3+}, so the numeral is the
        only thing distinguishing \ce{FeCl2} from \ce{FeCl3}.}
  \item Why are prefixes used in \ce{N2O4} but not in \ce{Al2O3}?
        \answerlines[2]{\ce{N2O4} is covalent, and the subscripts are not
        fixed by charge --- \ce{NO2} also exists, so prefixes are needed to
        say which. \ce{Al2O3} is ionic and its subscripts follow
        necessarily from \ce{Al^3+} and \ce{O^2-}, so prefixes would be
        redundant.}
  \item \ce{HClO3} and \ce{HClO2} differ by one oxygen. Describe how both
        the anion name and the acid name change.
        \answerlines[2]{The anion drops from chlorate to chlorite, and the
        acid follows it from chloric to chlorous --- \emph{-ate} pairs with
        \emph{-ic} and \emph{-ite} with \emph{-ous}.}
  \item A student writes \ce{CO} as ``carbon monooxide''. Correct the
        spelling and state the rule.
        \answerlines[2]{Carbon monoxide. When a prefix ending in a vowel
        meets an element name beginning with a vowel, the prefix's final
        vowel is dropped --- likewise pentoxide and heptoxide.}
\end{parts}

\begin{spiralstrip}{Chapter 2 \textbullet\ blocks 1--3}
  \item Protons in \ce{^{56}Fe^3+}: \answerblank[2cm]{26}
  \item Electrons in \ce{^{56}Fe^3+}: \answerblank[2cm]{23}
  \item Charge predicted for a group 5A monatomic ion:
        \answerblank[2cm]{$3-$}
  \item \ce{ClO3-} is named: \answerblank[2.6cm]{chlorate}
  \item An \emph{-ate} anion gives which acid ending?
        \answerblank[2cm]{\emph{-ic}}
\end{spiralstrip}

\end{document}
